\documentclass[12pt]{article}
\usepackage[margin=1in]{geometry}
\usepackage{enumitem}
\usepackage{titlesec}
\usepackage{parskip}
\usepackage{amsmath}
\usepackage{hyperref}
\usepackage{fontenc}

\title{\textbf{Kalyvas 2006 Claude Guide}\\
\large Kalyvas, Stathis N. \textit{The Logic of Violence in Civil War}.\\
New York: Cambridge University Press, 2006.}
\author{}
\date{}

\begin{document}

\maketitle
\tableofcontents
\newpage

%--------------------------------------------------
\section{Central Claims}
%--------------------------------------------------

\begin{itemize}[leftmargin=*, itemsep=4pt]
    \item Violence in civil war is not a direct, unmediated expression of the war's ``master cleavage''---the overarching ideological, ethnic, religious, or class division that supposedly defines the conflict (e.g., communist vs.\ nationalist, Hutu vs.\ Tutsi). Kalyvas calls the assumption that violence simply reads off this master cleavage the \textbf{``master cleavage fallacy''}---the central target of the book's critique of the existing civil war and genocide literatures.
    \item Instead, violence is a \textbf{joint product} of two distinct logics operating simultaneously: (1) the strategic/political logic of central actors (armies, insurgent organizations, political factions) competing for territorial control and civilian compliance, and (2) the private, local logic of individuals and communities pursuing personal disputes, feuds, land conflicts, and status rivalries that predate and outlast the war. Civil war provides an opportunity structure that allows local actors to ``piggyback'' their private agendas onto the war's public, ideological narrative---denouncing a neighbor as ``enemy'' or ``collaborator'' to settle a score over land, inheritance, or a marriage dispute.
    \item This joint production means violence is best understood as the outcome of an \textbf{alliance} between supralocal (political/military) actors, who need local information and manpower to identify and control the population, and local actors, who need the coercive power of the supralocal actor to prevail in local conflicts they could not otherwise win. Neither actor could produce the observed violence alone.
    \item Kalyvas draws a sharp analytical distinction between \textbf{indiscriminate violence} (applied to individuals based on group membership or category, without regard to individual behavior) and \textbf{selective violence} (targeted at specific individuals based on information about their actual or alleged actions/behavior). Selective violence requires \textbf{information}, which armed actors typically lack directly and must obtain from local civilian informants/denouncers.
    \item The type and intensity of violence used by an armed actor is a function of the actor's level of \textbf{territorial control}, not primarily of ideology, ethnicity, or regime type. Indiscriminate violence is a blunt substitute for information; it is used where control (and thus information) is low, precisely because the actor cannot identify individual defectors and instead punishes the population collectively as a deterrent. Selective violence requires enough control to gather intelligence (via denunciation) and to protect informants from retaliation, but not so much control that no one dares defect in the first place.
    \item Both indiscriminate and selective violence are therefore predicted to be \textbf{nonmonotonic} in control: violence (especially selective violence) peaks not where control is total or absent, but in \textbf{zones of contested or fragmented control}, where defection is possible, information is contestable, and the incumbent actor has enough presence to act on denunciations but not enough to deter defection outright.
    \item Identity and alliances in civil war are substantially \textbf{endogenous} to the conflict rather than fixed, primordial antecedents of it. Who counts as ``the enemy,'' who joins which side, and how cleavages are locally interpreted are shaped by the dynamics of the war itself (opportunism, coercion, protection-seeking, local rivalries)---not simply imported wholesale from a pre-existing ethnic or ideological cleavage.
    \item Indiscriminate violence is not only morally troubling but \textbf{strategically counterproductive} in the long run: because it punishes people regardless of individual behavior, it destroys the incentive to comply or collaborate (since compliance offers no protection), often driving the population toward the rival side. Selective violence, precisely because it is discriminating, preserves incentives for compliance and is therefore more effective at inducing civilian cooperation, though it depends on a fragile informational infrastructure of denunciation.
\end{itemize}

%--------------------------------------------------
\section{The Zones of Control Typology}
%--------------------------------------------------

\begin{itemize}[leftmargin=*, itemsep=4pt]
    \item Kalyvas develops a five-zone typology of territorial control, ranging from total control by actor A (Zone 1), through dominant A control with some incursion (Zone 2), contested/parity control (Zone 3), dominant B control (Zone 4), to total B control (Zone 5)---symmetric around the midpoint of contestation.
    \item \textbf{Zones of total control (1 and 5)}: The controlling actor has complete information (via a fully developed denunciation network, since defectors have nowhere to hide and collaborators need not fear reprisal) and no need for indiscriminate violence---overt violence tends to be low because compliance is already secured and deviants can be selectively identified and removed at low cost. Terror here is more selective, surgical, and rare, since the population is thoroughly incorporated.
    \item \textbf{Zones of dominant but incomplete control (2 and 4)}: The dominant actor has partial information; selective violence is feasible and increasingly used as the actor consolidates, since it has enough presence to protect informants and act on denunciations without yet facing total defection risk. This is where selective violence is often predicted to be most prevalent.
    \item \textbf{The zone of parity/contestation (Zone 3)}: Neither actor controls the area; sovereignty is fragmented and fluctuating. Information is hardest to obtain and protect (informants face high retaliation risk from the other side), so actors are more likely to fall back on indiscriminate violence as a blunt substitute for missing intelligence. Paradoxically, this is often the most violent zone overall, combining incidents from both indiscriminate and (attempted, if fragile) selective violence.
    \item The typology yields a core empirical prediction that Kalyvas tests directly against Greek village-level data: violence is not simply increasing or decreasing in control, but exhibits a more complex, actor-specific and often nonmonotonic pattern across the five zones---directly contradicting simpler theories that predict violence is monotonically decreasing (or increasing) in control, or is a simple function of military strength.
\end{itemize}

%--------------------------------------------------
\section{Denunciation and the Informational Basis of Selective Violence}
%--------------------------------------------------

\begin{itemize}[leftmargin=*, itemsep=4pt]
    \item Selective violence hinges on \textbf{denunciation}: local civilians providing information to armed actors about who among their neighbors is a genuine or alleged collaborator, sympathizer, or defector. Armed actors, especially outside occupying forces, typically lack the local knowledge to distinguish loyal civilians from opponents without this civilian intermediation.
    \item Denunciation is rarely a pure expression of ideological conviction. Kalyvas shows empirically (using Greek Civil War archives) that denunciations frequently encode \textbf{private motives}---land disputes, inheritance conflicts, sexual/marital jealousy, intra-family feuds, status competition, and long-standing local rivalries that predate the civil war entirely. The war supplies a vocabulary (``collaborator,'' ``traitor,'' ``fascist,'' ``communist'') and an enforcement mechanism (the armed actor) that individuals exploit to settle these private scores.
    \item This generates the core micro-mechanism of the book: local actors trade information for protection/power, and supralocal actors trade coercive capacity for local intelligence and legitimacy---an \textbf{exchange relationship} that binds the ``public'' war to a substrate of ``private'' conflicts.
    \item Because denunciation exposes the informant to retaliation risk (from the targeted individual's family/allies, or from the rival armed actor if control shifts), the volume and reliability of denunciation is itself a function of the protective capacity of the controlling actor---linking back to the zones-of-control argument: only where control is stable enough to shield informants does a functioning denunciation network emerge.
    \item Kalyvas also stresses that denunciation is a two-edged, self-reinforcing process: once a village or family has been targeted, retaliatory denunciation in the opposite direction becomes likely if control later shifts, producing cycles of tit-for-tat local violence that persist independent of the war's broader strategic trajectory.
\end{itemize}

%--------------------------------------------------
\section{The Master Narrative vs.\ Private/Local Conflicts}
%--------------------------------------------------

\begin{itemize}[leftmargin=*, itemsep=4pt]
    \item Kalyvas draws a foundational distinction between the \textbf{``master narrative''} of the civil war---the public, national-level story of ideological or ethnic confrontation that dominates elite discourse, historiography, and after-the-fact collective memory---and the actual, observed \textbf{local dynamics of violence}, which are frequently driven by concerns only loosely (or not at all) connected to that master narrative.
    \item This distinction reframes the analytical unit of civil war violence: rather than treating ``the war'' as a single coherent process to be explained by macro-level variables (regime type, ideology, foreign intervention), Kalyvas insists that civil war violence is fundamentally a \textbf{local, transaction-level phenomenon} best studied at the level of the village, the family, or even the individual denunciation.
    \item Local ``private'' conflicts are not simply epiphenomenal noise around the ``real'' war; they are constitutive of how the war is actually fought and experienced on the ground. The war's public ideological cleavage supplies the \emph{institutional technology} (armed organizations, categories of enmity, mechanisms of punishment) that local actors then capture and repurpose for private ends.
    \item This has a major methodological implication: aggregate, national-level, or purely ideological/ethnic accounts of civil war violence \textbf{systematically mismeasure} its causes, because they attribute to ``the war'' (the master cleavage) violence that is actually produced by local dynamics that happen to be legible only through the vocabulary the war supplies.
\end{itemize}

%--------------------------------------------------
\section{Endogenous Identity and Alliance Formation}
%--------------------------------------------------

\begin{itemize}[leftmargin=*, itemsep=4pt]
    \item Rather than treating wartime identities and alliances (who is ``with'' the guerrillas, who is ``with'' the government) as fixed attributes that individuals bring into the conflict, Kalyvas treats them as largely \textbf{endogenous outcomes} of the war's unfolding dynamics---shaped by coercion, opportunism, the search for protection, and the shifting fortunes of local control.
    \item Individuals and villages frequently change sides, or perform allegiance to whichever actor currently controls their area (a survival strategy), meaning observed ``support'' for an armed actor is frequently a function of \textbf{control}, not of prior conviction---an important corrective to studies that infer ideological commitment directly from wartime behavior.
    \item This endogeneity argument parallels and extends a broader critique (also present in ethnic conflict studies) that treating identity as a fixed, exogenous variable that predicts violence gets the causal arrow partly backwards: violence and the war itself help construct and harden the very cleavages they are often said to merely express.
\end{itemize}

%--------------------------------------------------
\section{Empirical Strategy: The Greek Civil War and Comparative Cases}
%--------------------------------------------------

\begin{itemize}[leftmargin=*, itemsep=4pt]
    \item The empirical core of the book is a micro-level, village- and district-level study of violence during the \textbf{Greek Civil War} (and the preceding Axis Occupation, 1941--1944), drawing on local archives, judicial records, oral histories, and village-level data on control, denunciations, and killings.
    \item Kalyvas assembles a novel dataset coding \textbf{territorial control} and \textbf{violence incidents} at the village level over time, allowing him to test the zones-of-control predictions against actual patterns of selective and indiscriminate killing---a level of empirical granularity rare in civil war scholarship, which typically relies on national- or regional-level aggregates.
    \item The choice of Greece is deliberate: it offers rich local documentation, a civil war widely (and, per Kalyvas, wrongly) understood through a simple left/right ideological master narrative, and sufficient variation in control and violence across villages and time to identify the proposed mechanisms.
    \item Kalyvas supplements the Greek case with comparative evidence and illustrations from other conflicts---including the Vietnam War, the Spanish Civil War, the Algerian War, and various post-colonial and contemporary civil wars---to argue that the joint-production logic and the zones-of-control predictions are not idiosyncratic to Greece but generalize across highly disparate conflicts.
    \item Methodologically, the book models an explicit \textbf{multi-level research design}: national/regional narrative and control data contextualize village-level violence data, and qualitative process-tracing of specific denunciation episodes (drawn from archival and interview material) supplies the micro-mechanism evidence that a purely quantitative cross-village analysis could not establish alone.
\end{itemize}

%--------------------------------------------------
\section{Critique of Collective/Categorical Approaches to Violence}
%--------------------------------------------------

\begin{itemize}[leftmargin=*, itemsep=4pt]
    \item Kalyvas directly challenges collective or \textbf{categorical} approaches to mass violence---including much of the genocide studies literature---that assume violence is targeted uniformly at members of a group based solely on group membership (ethnicity, class, religion), driven by top-down ideological intent.
    \item He argues that even in cases conventionally coded as ``ethnic cleansing'' or ``genocide,'' close empirical examination typically reveals substantial \textbf{selectivity and local variation}: perpetrators often rely on local informants to identify specific individuals, targeting is uneven across villages and families in ways that categorical group-based accounts cannot explain, and much of the ``who gets killed'' variation tracks local control and denunciation dynamics rather than uniform application of a genocidal category.
    \item This critique does not deny that group categories matter---they clearly shape who is eligible to be targeted and supply the discourse of enmity---but insists that categorical membership is a necessary, not sufficient, condition for victimization; \textbf{actual} targeting decisions are mediated by local informational and strategic dynamics that categorical theories omit.
    \item The broader implication is methodological: scholars of mass violence should disaggregate their unit of analysis below the national or ethnic-group level, examining micro-level variation in violence across localities, families, and individuals, rather than treating ``the group'' as a sufficient explanatory category.
\end{itemize}

%--------------------------------------------------
\section{Methodological Contributions and Limitations}
%--------------------------------------------------

\begin{itemize}[leftmargin=*, itemsep=4pt]
    \item \textbf{The micro-comparative turn}: Kalyvas is a foundational text in the ``micro-dynamics of civil war'' research program (alongside scholars like Elisabeth Wood and Jeremy Weinstein), pushing the field away from country-year cross-national regressions of civil war onset/duration toward disaggregated, subnational, often village- or event-level analysis of violence.
    \item \textbf{Novel data construction}: The village-level control and violence dataset for Greece is a significant original empirical contribution, built from painstaking archival and interview work rather than relying on pre-existing cross-national datasets---a template later followed by much of the micro-level civil war literature.
    \item \textbf{Endogeneity and selection concerns}: Territorial control and violence are plausibly jointly determined (violence can itself shift control, and anticipated control can shape strategies of violence), raising causal identification challenges that Kalyvas addresses primarily through careful process-tracing and timing evidence rather than a fully worked-out identification strategy.
    \item \textbf{Generalizability from a single primary case}: While Kalyvas offers illustrative comparative evidence from other wars, the core empirical tests rest heavily on the Greek case; whether the zones-of-control predictions hold with the same functional form (especially the nonmonotonicity claim) in wars with very different technologies of violence (e.g., aerial bombardment, genocide with limited local intermediation, high-tech counterinsurgency) is left more open than the Greek results might suggest.
    \item \textbf{Rational-choice/microfoundational framing}: The theory is explicitly actor-centered and strategic (armed actors as goal-directed organizations optimizing control and information; civilians as strategic actors managing risk and opportunity), which gives the argument analytical precision and testable predictions, but has drawn criticism for downplaying ideational, emotional, and normative motivations for violence and collaboration that are harder to formalize.
    \item \textbf{Conceptual clarity as contribution}: Perhaps the book's most durable contribution is definitional/conceptual: crisply distinguishing selective from indiscriminate violence, and civil war ``master cleavage'' from local private conflict, gives the field an analytical vocabulary that has been widely adopted even by scholars who do not share Kalyvas's specific zones-of-control predictions.
\end{itemize}

%--------------------------------------------------
\section{Connections to the Broader Literature}
%--------------------------------------------------

\begin{itemize}[leftmargin=*, itemsep=4pt]
    \item \textbf{Huntington 1968}: Huntington's core concern is the gap between political mobilization/participation and institutionalization, with violence as a symptom of weak institutions unable to channel mobilized demands. Kalyvas offers a complementary but more granular claim: even within a single ``institutionalization failure'' (civil war), the specific pattern of violence is governed by local territorial control and information, not simply by the aggregate weakness of central institutions. Both share a skepticism of treating violence as a direct, unmediated output of high-level structural conditions (modernization for Huntington, ideology/ethnicity for the civil war literature Kalyvas critiques).
    \item \textbf{Dahl 1972 (Polyarchy)}: Dahl's conditions for stable polyarchy include low coerciveness of the regime and mutual security among contending groups. Kalyvas's zones-of-control analysis can be read as a micro-level account of exactly what happens when mutual security collapses locally---contested-control zones (Zone 3) are precisely where the guarantees Dahl treats as prerequisites for peaceful contestation are absent, producing violence as a substitute for institutionalized competition.
    \item \textbf{Moore 1966 (Social Origins of Dictatorship and Democracy)}: Moore's macro-historical, class-coalition account of revolutionary and reactionary violence operates at a level of aggregation (national class coalitions, agrarian structures) that Kalyvas's book implicitly critiques: Moore's framework would predict violence patterns from class/ideological cleavage, exactly the master-cleavage inference Kalyvas shows to be empirically unreliable at the local level. The two works are usefully read against each other as macro-structural vs.\ micro-strategic accounts of the same broad phenomenon (violent political conflict).
    \item \textbf{Putnam 1993}: Putnam's civic community/social capital argument and Kalyvas's denunciation dynamics both hinge on the density and structure of local social networks, but with opposite valence: in Putnam, dense horizontal networks and trust generate cooperative civic outcomes; in Kalyvas, dense local networks of information (who knows what about whom) become the infrastructure that enables selective violence when captured by armed actors. Both works insist that national-level political outcomes are mediated by highly local social structure---a shared commitment to disaggregating the unit of analysis below the nation-state.
    \item \textbf{Cox 1997 (Making Votes Count)}: While Cox's strategic-coordination logic concerns electoral competition rather than violence, both Cox and Kalyvas model political actors as strategically adapting behavior to institutional/structural incentives (electoral rules for Cox; zones of territorial control for Kalyvas), and both derive sharp, falsifiable comparative-statics predictions from a spare rational-actor model---an affinity in methodological style across substantively distant subfields.
    \item \textbf{Lijphart 1999 (Patterns of Democracy)}: Lijphart's consensus/majoritarian typology addresses how divided societies manage conflict through institutional design once violence is off the table; Kalyvas's book can be read as addressing the prior, darker question of what happens when institutional accommodation fails or has not yet been built, and cleavages are instead processed through coercion rather than power-sharing. Both scholars resist essentialist readings of ethnic/ideological divisions as automatically producing particular political outcomes (violence for Kalyvas, instability for naive readings of Lijphart's plural societies).
    \item \textbf{Weber 1930/1968}: Kalyvas's zones-of-control argument is a direct, empirically elaborated application of the Weberian premise that authority requires an effective monopoly of legitimate violence over a territory; where that monopoly is only partial (the contested Zone 3), neither actor can achieve Weberian domination, and the resulting order relies on raw coercion and information extraction rather than legitimacy---an illustration of what politics looks like in the absence of a stabilized Weberian state.
    \item \textbf{Huber and Shipan 2002 (Deliberate Discretion)}: Both works share a concern with principal-agent problems and information asymmetries in governance---Huber and Shipan model legislatures delegating discretion to bureaucratic agents under incomplete information, while Kalyvas models armed actors relying on local civilian ``agents'' (informants/denouncers) to overcome their own severe informational deficits about the population. In both, the controlling actor's institutional design (statutory specificity; investment in a denunciation apparatus) is a direct response to the underlying information problem.
    \item \textbf{Powell 2000 (Elections as Instruments of Democracy)}: Powell's concern with the chain of representation and accountability presupposes exactly the kind of stable, peaceful political order that Kalyvas's civil war context lacks; Kalyvas can be read as illuminating the pathological substitute for representation and accountability that emerges when normal political competition breaks down---denunciation and violence become informal (and terroristic) functional analogues to the electoral sanctioning mechanisms Powell studies.
    \item \textbf{Stokes 2013 (Brokers, Voters, and Clientelism)}: Stokes's account of clientelistic brokers who use fine-grained local information to target voters with selective incentives is structurally analogous to Kalyvas's account of denunciation networks that use fine-grained local information to target individuals with selective violence. Both are, at core, theories of how local intermediaries with private information enable a supralocal principal (party machine; armed actor) to act selectively rather than indiscriminately---clientelism and selective violence are, in this sense, close theoretical cousins operating through parallel principal-agent-broker structures.
    \item \textbf{Svolik 2012 (The Politics of Authoritarian Control)}: Svolik's analysis of authoritarian regimes' reliance on repression and the loyalty-competence tradeoffs among agents of coercion parallels Kalyvas's account of how armed actors must balance indiscriminate repression (blunt, low-information-cost, but corrosive to compliance) against selective repression (high-information-cost but more effective and legitimacy-preserving). Both authors treat repression/violence as a costly, information-constrained strategic choice rather than an unconstrained expression of a regime's or faction's preferences.
    \item \textbf{Carey 2009 (Legislative Voting and Accountability)}: A more distant connection, but both works are attentive to how monitoring and information (voter/legislator behavior for Carey; civilian collaboration/defection for Kalyvas) discipline agents' behavior, and both emphasize that the \emph{capacity to observe and sanction} individual behavior (rather than the mere existence of formal rules or professed cleavages) determines political outcomes.
\end{itemize}

\end{document}
